In the previous section we proved our main theorem, Theorem~\ref{subsurface theorem}, except for two key lemmas.  We state and prove those lemmas in this section.

Lemmas~\ref{main lemma easy} and~\ref{main lemma corollary} apply to Case~1 and Case~2 of our proof of Theorem~\ref{subsurface theorem}, respectively.  Lemma~\ref{main lemma easy} applies to the $\dd$-disc case, and Lemma~\ref{main lemma corollary} applies to the $\dd$-half-disc case.  Lemma~\ref{main lemma easy} is easy to prove, and Lemma~\ref{main lemma corollary} is harder.  We present Lemma~\ref{main lemma corollary} as a corollary of Lemma~\ref{main lemma}, the Main Lemma.  We introduce several definitions before proceeding to the proofs.

Let $H$ be a handlebody.  Let $S$ be a compact surface in $\partial H$.  Let $\dd$ be a disc system in $H$.  Assume that $\dd$ is taut with respect to $\partial S$.  We defined taut in Section~\ref{maintheorem}.  Also in that section, we defined the map $\m$ associated to the triple $(H, S, \dd)$.  The map $\m$ is a spherical diagram showing the roads $R$ and the collection $\dpm$ of countries.  A $\dd$-interface $\dint{a}$ for $\dpm$ colors each country in $\m$ either black, white, or gray.

In Section~\ref{maintheorem}, we explained how to view $\m$ as a disc by choosing one country to be the outside.  This lets us refer to the inside and outside of simple closed curves in $\m$.  In this section, we will always view maps $\m$ in this way.  We denote the boundary of the outside country $\bd\m$.  When $\dint{a}$ is a $\dd$-half-disc, we choose the gray disc to be the outside country.  Otherwise, our convention is to pick a black disc as the outside.

Every road $R$ in $\m$ has $p$ \emph{prongs}, where $2p$ is the number of sides of $R$.  For example, the roads in Figure~\ref{segregation} have five, four, and four prongs, respectively.  Each road in $\m$ has at least two prongs.  Define the Euler characteristic $\chi(R)$ of a road $R$ to be $\frac{2-p}{2}$.  For example, a rectangular road has zero Euler characteristic.  All other roads in $\m$ have negative Euler characteristic.  Define the Euler characteristic of a collection of roads to be the sum of the Euler characteristics of the roads.

Observe that the Euler characteristic of the collection of all roads in $\m$ equals $\chi(S)$.  This can be seen by considering the graph dual to $S$.  View the roads inside $S$ rather than inside $\m$.  Put a vertex inside each road of $S$.  For each place where two roads border, connect the corresponding vertices with an edge.  See Figure~\ref{dual graph} for an example.  The dual graph has the same Euler characteristic as the surface $S$ because the dual graph is a deformation retract of $S$.  The contribution near a vertex to the Euler characteristic of the graph equals the contribution of the corresponding road to $\chi(S)$.

\begin{figure}
% dual graph
\epsfig{file=figures/dual_graph.eps, height=1.5in}
\caption{
Each 3-pronged road in this picture contributes $-1/2$ to the Euler characteristic of $S$, for a total Euler characteristic of $-1$.  The dual graph also has Euler characteristic $-1$.  Each 3-pronged vertex contributes $-1/2$ to the Euler characteristic of the graph.
}
\label{dual graph}
\end{figure}

Let $R$ be a road of $\m$.  Define $d(R)$ to be $\min(1,\ec{R})$.  Then $d(R)$ is $0$, $1/2$, or $1$ depending on whether $R$ has two, three, or more than three prongs.  Let $X$ be a subset of $\m$.  Recall the definition of a non-segregating road in Section~\ref{maintheorem}.  Define $d(X)$ to be the sum of $d(R)$ over all \emph{non-segregating} roads $R$ that intersect $X$.  Loosely speaking, the function $d$ counts the non-segregating roads of negative curvature that intersect the subset $X$ of $\m$.  A non-segregating road $R$ with $d(R)$ equal $1/2$ must be a 3-pronged road touching a gray $\bd\m$.  Observe that any subset $X$ of $\m$ satisfies $d(X)\leq\ec{S}$.

We now define the continents of $\m$.  Consider the union in $\m$ of (1) all the countries in $\m$, (2) the boundary circle $\partial \m$ if it corresponds to a black country, and (3) all the arcs of $\partial S$ that do not connect a black country to a white country.  Let $U$ be a regular neighborhood of this subset in $\m$.  See the left sides of Figures~\ref{map1} and~\ref{graymap} for examples.  Each component of $\partial U$ is either a circle or an arc.  Arc components exist only when $\dint{a}$ is a $\dd$-half-disc (see Figure~\ref{graymap}).  In this case, all arcs begin and end at the boundary circle $\partial \m$.

\begin{figure}
% map showing a D-DISC 
\epsfig{file=figures/blackmap1.eps, height=2.25in}%
\hspace{.5in}%
\epsfig{file=figures/blackmap2.eps, height=2.25in}
\caption{
We show two pictures of the same map $\m$ with a $\dd$-disc.  The outer country is black.  We show the collection $U$ of continents in gray.  The left picture shows $U$ before compressing.  The right shows $U$ after compressing.  There is a black continent and a white continent.
}
\label{map1}
\end{figure}

Compress the inessential boundary components of $U$
as follows.
Fill in each circle of $\partial U$ that bounds a disc in $\m-\dd$.
Boundary-compress every arc of $\partial U$ that is inessential
in $\m-\dd$.
Each component of $\bd U$ is now either an essential circle in 
$\m-\dd$ or an essential arc in $\m-\dd$ based at $\partial \m$.  
See the right sides of Figures~\ref{map1} and~\ref{graymap} for
examples.

\begin{figure}
% 
\epsfig{file=figures/graymap1.eps, height=2.25in}
\hspace{.5in}
\epsfig{file=figures/graymap2.eps, height=2.25in}
\caption{
We show a map $\m$ with a $\dd$-half-disc.  The boundary $\bd\m$ corresponds to a gray country.  The left shows the collection $U$ of continents before compressing, and the right after compressing.  Again, there is a black continent and a white continent.}
\label{graymap}
\end{figure}

A \emph{continent} of $\m$ is a connected component of the
neighborhood $U$ after compressing.  Each continent 
contains at least one country.
Each continent contains only like-colored countries.
Thus we can speak of black and white continents.

\begin{lemma}
\label{main lemma fact}
Let $U$ be the collection of continents for a $\dd$-interface
$\dint{a}$, as defined above.  Let $m$ be the number of 
segregating roads.  
Let $K$ be a component of $\partial U$ that intersects each road 
in at most one arc.  Then $|K\cap \partial S|\leq 2d(K)+2m$.
\end{lemma}

\begin{proof}
The curve $K$ intersects the boundary of a non-segregating road with 
two prongs in zero points, with three prongs in at most one point, 
and with four or more prongs in at most two points.
The curve $K$ also intersects every segregated road in at most 
two points.  The result follows.
\end{proof}

We now define a partial order $\subseteq$ on the components of 
$\partial U$.  Let $K_1$ and $K_2$ be distinct components of 
$\partial U$.  We say $K_1\subseteq K_2$ if there is a
simple arc $k$ in some road so that $K_1$ is inside $K_2\cup k$
(see Figure~\ref{partialorder}).
We also define $K_1\subseteq K_1$ for all components $K_1$
of $\bd U$.  This makes $\subseteq$ a partial order on the components
of $\partial U$.

\begin{figure}
%
\epsfig{file=figures/partialorder.eps, height=2.25in}%
\caption{
These two components $K_1$ and $K_2$ of $\bd U$ satisfy 
$K_1\subseteq K_2$.  The curve $K_1$ lies inside the union
of $K_2$ and the dotted arc.
}
\label{partialorder}
\end{figure}

Observe that if $K_1\subseteq K_2$ for distinct $K_1$ and $K_2$, then $K_1$ is a circle rather than an arc.  Let $K_2$ be a component of $\partial U$.  Observe that there is a distinct component $K_1$ of $\partial U$ such that $K_1\subseteq K_2$ if and only if $K_2$ intersects a road in more than one arc.  In particular, the boundary $\partial U$ contains a circle.

We now state Lemma~\ref{main lemma easy}.
The proof we give specializes our proof of the Main Lemma~\ref{main lemma}.

\begin{lemma}[cf. Lemma~\ref{main lemma corollary}]
\label{main lemma easy}
Let $H$ be a handlebody, and let $S$ be a compact surface in 
$\partial H$.  Let $\dd$ be a taut disc system with respect 
to $\partial S$.  Let $\dint{b}$ be a $\dd$-disc with 
$n$ segregated roads.
Then there is an essential disc in $H-\dd$ that intersects $S$ in no more
than $\ec{S}+n$ arcs.
\end{lemma}

\begin{proof}
Consider the map $\m$ corresponding to $(H,S,\dd)$, with 
$\dd$-disc $\dint{b}$.
Let $U$ denote the union of continents in $\m$, as defined 
above.  Since no countries are gray, all components of 
$\partial U$ are circles.  Let $C$ be a circle that is minimal 
with respect to the partial ordering $\subseteq$.  Then $C$ 
intersects each road in at most one arc.
Apply Lemma~\ref{main lemma fact} to $C$.  Then 
$|C\cap\partial S|\leq 2d(C)+2n\leq 2\ec{S}+2n$.  Thus
$C$ intersects $S$ in at most $\ec{S}+n$ arcs.  Since $C$
is disjoint from $\dd$ in $H$, it bounds a disc in $H$.
\end{proof}

We now state and prove the Main Lemma.

\begin{lemma}[\textbf{Main Lemma}]
\label{main lemma}
Let $H$ be a handlebody, and let $S$ be a compact surface in $\partial H$.  Let $\dd$ be a taut disc system with respect to $\partial S$.  Let $\dint{c}$ be a $\dd$-half-disc with $m$ segregating roads. Then in $H-\dd$ there is either (1) an essential circle $C$ such that $|C\cap\partial S|\leq2(\ec{S}+2m+1)$, or (2) a half-disc $K$ such that $|K\cap\partial S|\leq\ec{S}+2m+1$.
\end{lemma}

\begin{proof}
Let $\m$ be the associated map.  The $\dd$-interface $\dint{c}$ gives a color to each disc in $\dpm$.  Since $\dint{c}$ is a $\dd$-half-disc, exactly one disc in $\dpm$ is gray.  The discs in $\dpm$ correspond to the countries of $\m$.  Following our convention, the disc corresponding to $\bd \m$ is the gray disc.

Let $U$ denote the union of continents in $\m$.  Assume temporarily that $\partial U$ contains a circle.  Let $C$ be a circle in $\bd U$ that is minimal with respect to the partial ordering $\subseteq$.  Then $C$ intersects each road in at most one arc.  Recall this by examining Figure~\ref{partialorder}.  Applying Lemma~\ref{main lemma fact}, we have $|C\cap\partial S|\leq 2d(C)+2m\leq 2\ec{S}+2m\leq 2(\ec{S}+2m+1)$.  We are done in this case, so we can assume otherwise.

\textbf{Assume that every component of $\bd U$ is an arc.}  Then every component $K$ of $\bd U$ intersects each road in at most one arc.  We will construct an arc $K$ of $\bd U$ that satisfies the conclusion of the lemma.  

Give the complement of $U$ in the disc $\m$ the structure of a graph, as follows.  The disc $\m$ contains a collection of arcs coming from $\bd S$.  These arcs bound the roads in $\m-\dpm$.  Let each component of the complement of these arcs in $\m-U$ be a vertex.  Note that some of these components correspond to roads, and others do not.  Let each arc in $\m-U$ be an edge.  The incidence relation gives these vertices and edges the structure of a finite graph.  Observe that each road in $\m$ corresponds to only one vertex in the graph.  Note also that this graph may be disconnected.

Let $\Gamma$ be a connected component of this graph.  Since $\bd U$ contained no circles, the graph $\Gamma$ is a tree.  Observe that every second vertex of $\Gamma$ corresponds to a road of $\m$.  More precisely, let $v$ be a vertex of $\Gamma$ that corresponds to a road.  The vertices of $\Gamma$ that correspond to roads are exactly the vertices that are an even number of edges away from $v$.

%  We are now going to backtrack and modify the collection $U$ of continents in $\m$ slightly.  We will then perform the same graph construction on the new collection $U$.  The effect on $\Gamma$ may be to divide $\Gamma$ into several smaller graphs.

Define the \emph{extremal} vertices of a tree to be the vertices of the tree that touch only one edge.  Define the \emph{interior} vertices of a tree to be the non-extremal vertices.  Observe that if an extremal vertex of the tree $\Gamma$ corresponds to a road, then that vertex corresponds to a road touching $\bd\m$ with three or more prongs.  Note also that roads touching $\bd\m$ are non-segregating. 

Assume temporarily that no interior vertex of $\Gamma$ corresponds to a non-segregating road.  Then every interior vertex that corresponds to a road corresponds to a segregating road.  Choose a geodesic path $P$ in $\Gamma$ that connects two extremal vertices.  Realize $P$ as a half-disc $K$ in $\m$ based at the gray disc corresponding to $\bd\m$.  Observe that $|K\cap S|$ equals the number of edges of $P$.  There are two cases.

\textbf{Case 1: Neither endpoint of $P$ corresponds to a road.}  Then the number of edges of $P$ is less than or equal to twice the number of segregating roads in $\m$.  Hence $|K\cap S|\leq 2m<\ec{S}+2m+1$, and we are done with this case. 

\textbf{Case 2. At least one endpoint of $P$ corresponds to a road.}  Then at least one road in $\m$ has three or more prongs, so $\ec{S}>0$, or equivalently $\ec{S}\geq 1$.  The number of edges of $P$ is less than or equal to $2m+2$, since each endpoint of $P$ can contribute an additional edge.   Hence $|K\cap S|\leq 2m+2\leq\ec{S}+2m+1$.  We are done with the second case and with our temporary assumption. 

\textbf{Assume that at least one interior vertex of $\Gamma$ corresponds to a non-segregating road.}  Such an interior vertex must correspond to a road with four or more prongs, so $\ec{S}\geq 1$.

\begin{figure}
% dual graph
\epsfig{file=figures/four_paths.eps, height=2in}
\caption{
We illustrate how four paths can be chosen in the graph $\Gamma$ to satisfy the desired properties.  The subgraph $\Gamma_{0}$ is shown in bold, though much of the graph $\Gamma$ is not shown.  Each of the four paths intersects exactly one extremal vertex of $\Gamma_{0}$.  No more than two of the four paths intersect any one extremal vertex of $\Gamma$.
}
\label{four paths}
\end{figure}

Let $\Gamma_{0}$ denote the convex hull in $\Gamma$ of the interior vertices that correspond to non-segregating roads.  By assumption this is nonempty.  Choose four geodesic paths $P_{1}$, $P_{2}$, $P_{3}$, and $P_{4}$ in $\Gamma$ such that (1) no more than two of the $P_{i}$ touch any extremal vertex of $\Gamma$, and (2) each $P_{i}$ intersects $\Gamma_{0}$ in at most one vertex.  This can be done easily.  The process is illustrated in Figure~\ref{four paths}.

The argument now divides into cases according to how many extremal vertices of $\Gamma$ correspond to roads.  The argument is similar to above.  Note that if distinct extremal vertices of $\Gamma$ correspond to roads, they correspond to distinct roads.

Represent each geodesic path $P_{i}$ in $\Gamma$ by a half-disc $K_{i}$ in $\m$ based at the gray disc $\bd\m$.  Observe that $|K_{i}\cap S|$ equals the number of edges of $P_{i}$.  At most four edges of each $P_{i}$ do not correspond to the intersection of $K_{i}$ with the boundary of a segregating road: possibly two edges for an extremal vertex of $\Gamma_{0}$ and possibly one edge for each extremal vertex of $P_{i}$.  There are three cases.

\textbf{Case 1. Zero or one extremal vertices of $\Gamma$ correspond to roads.}
Choose one of the four paths, call it $P_{1}$, that does not intersect an extremal vertex of $\Gamma$ that corresponds to a road.  Then at most two edges of $P_{1}$ correspond to non-segregating roads.  The rest correspond to segregating roads.  Thus we have $|K_{1}\cap S|\leq 2m+2 \leq \ec{S}+2m+1$.  The latter inequality holds since $\ec{S}\geq 1$.

\textbf{Case 2. Two or three extremal vertices of $\Gamma$ correspond to roads.}
Then at least one of the four paths, call it $P_{1}$, intersects at most one extremal vertex of $\Gamma$ that corresponds to a road.  Such a road must be non-segregating and contributes one additional edge.  Moreover, each extremal vertex contributes at least $1/2$ to $\ec{S}$, so that $\ec{S}\geq 2$.  Thus $|K_{1}\cap S|\leq 2m+3 \leq \ec{S}+2m+1$.

\textbf{Case 3. Four or more extremal vertices of $\Gamma$ correspond to roads.}
Then each of the four paths intersect at most two extremal vertices of $\Gamma$ corresponding to a road.  This contributes at most two additional edges.  Moreover, each extremal vertex contributes at least $1/2$ to $\ec{S}$, so that $\ec{S}\geq 3$.  Taking $P_{1}$ for example, we have $|K_{1}\cap S|\leq 2m+4 \leq \ec{S}+2m+1$.

In each case we found a half-disc $K$ such that $|K\cap S|\leq \ec{S}+2m+1$, so we are done.
\end{proof}

\begin{lemma}[corollary of Main Lemma~\ref{main lemma}; 
cf. Lemma~\ref{main lemma easy}]
\label{main lemma corollary}
Let $H$ be a handlebody, and let $S$ be a compact surface in $\partial H$.
Let $\dd$ be a locally minimal disc system with respect to $\partial S$.
Let $\dint{c}$ be a $\dd$-half-disc with $m$ segregated roads.
Then there is an essential disc in $H-\dd$ that intersects $S$ in no more than
$\ec{S}+2m+1$ arcs.
\end{lemma}

\begin{proof}
Let $D$ be the disc in $\dpm$ at which the $\dd$-half-disc $\dint{c}$ is based.  Apply Lemma~\ref{main lemma} to $\dint{c}$.  There are two cases.

\textbf{Case 1.}  There is an essential circle $C$ in $\m-\dpm$ such 
that $|C\cap\partial S|\leq2\ec{S}|+4m+2$.  Then $C$ intersects
$S$ in no more than $\ec{S}+2m+1$ arcs.  The circle $C$ bounds
a disc in $H-\dd$, so we are done.

\textbf{Case 2.}There is a half-disc $K$ based at $D$ such that $|K\cap\partial S|\leq\ec{S}+2m+1$.  The endpoints $\partial K$ divide $\partial D$ into two components $D_B$ and $D_W$.  Since $\dd$ is locally minimal, the half-disc $K$ is not a short-cut
half-disc.  Therefore $D_B$, say, satisfies $|D_B\cap\partial S|\leq |K\cap\partial S|$.  The circle $C$ formed by $D_B$ and $K$ is essential and satisfies $|C\cap\partial S|\leq 2\ec{S}+4m+2$, so we can finish as in Case~1.
\end{proof}

